% GENERATED FILE. Edit canonical lessons, then run npm run generate:books.
% canonical-source-hash: fnv1a64:cfa934fe74df1e0c
% canonical-lessons: GE-C449-kassierer, GE-C449-kassieren, GE-C449-bargeld, GE-C449-kartenlesegeraet, GE-C449-wechselgeld

\chapter{An der Kasse}
\label{ch:ge-an-der-kasse}
\hlchaptermodality{449}

\begin{chapteropening}
\textbf{By the end of this chapter:} \emph{I can pay at the till: the cashier, taking payment, cash, the card reader and the change.}

Complete the last lesson of chapter 449: I can pay at the till: the cashier, taking payment, cash, the card reader and the change.
\end{chapteropening}

\section[der Kassierer]{der Kassierer — the cashier}

\label{lesson:GE-C449-kassierer}



\begin{quote}
Before the new one: say the German for the small parcel, then the German for the registered letter.
\end{quote}

\subsection*{You'll want to know: der Kassierer}
\textbf{der Kassierer} — \textquotedblleft{}the cashier\textquotedblright{}.

\begin{quote}
Der Kassierer gibt mir die Quittung.
\end{quote}

\begin{quote}
\emph{The cashier gives me the receipt.}
\end{quote}

Same form in the plural. A woman: \emph{die Kassiererin}.

\begin{cousinweb}
\emph{Kassierer} is built on \emph{kassieren} 'to take money in', from \emph{Kasse} 'till, cash box', which came from Italian \emph{cassa}, from Latin \emph{capsa} 'box' (EtymDB).
\end{cousinweb}

\begin{grammarlens}[title={What you've built}]
One more word for telling, explaining and describing at B1.
\end{grammarlens}

\subsection*{Your turn}
\begin{itemize}
\raggedright
  \item \emph{Say it:} \emph{der Kassierer}
  \item \emph{Say it:} \emph{der Kassierer}, once more
  \item \emph{Say it:} a sentence of your own with \emph{der Kassierer}
  \item \emph{Recall:} say the German for the small parcel, then the German for the registered letter, then say \emph{der Kassierer} again
\end{itemize}

\begin{tcolorbox}[breakable,colback=teal!4,colframe=teal!35!black,title={Before you move on}]
What does der Kassierer mean? (\textquotedblleft{}The cashier\textquotedblright{}.) Say it once more.
\end{tcolorbox}
\section[kassieren]{kassieren — to take the money, to collect (payment)}

\label{lesson:GE-C449-kassieren}



\begin{quote}
Before the new one: say the German for the registered letter, then the German for the cashier.
\end{quote}

\subsection*{You'll want to know: kassieren}
\textbf{kassieren} — \textquotedblleft{}to take the money, to collect (payment)\textquotedblright{}.

\begin{quote}
Im Restaurant kassiert der Kellner an jedem Tisch.
\end{quote}

\begin{quote}
\emph{In the restaurant the waiter takes the money at each table.}
\end{quote}

Past participle without \emph{ge-}: \emph{er hat kassiert}. Colloquial: \emph{eine Strafe kassieren} — to get a penalty.

\begin{cousinweb}
\emph{kassieren} is built on \emph{Kasse} 'till, cash box' (EtymDB lists the verb), which came from Italian \emph{cassa}, from Latin \emph{capsa} 'box' (Lewis \& Short).
\end{cousinweb}

\begin{grammarlens}[title={What you've built}]
One more word for telling, explaining and describing at B1.
\end{grammarlens}

\subsection*{Your turn}
\begin{itemize}
\raggedright
  \item \emph{Say it:} \emph{kassieren}
  \item \emph{Say it:} \emph{kassieren}, once more
  \item \emph{Say it:} a sentence of your own with \emph{kassieren}
  \item \emph{Recall:} say the German for the registered letter, then the German for the cashier, then say \emph{kassieren} again
\end{itemize}

\begin{tcolorbox}[breakable,colback=teal!4,colframe=teal!35!black,title={Before you move on}]
What does kassieren mean? (\textquotedblleft{}To take the money, to collect (payment)\textquotedblright{}.) Say it once more.
\end{tcolorbox}
\section[das Bargeld]{das Bargeld — the cash}

\label{lesson:GE-C449-bargeld}



\begin{quote}
Before the new one: say the German for the cashier, then the German for to take the money.
\end{quote}

\subsection*{You'll want to know: das Bargeld}
\textbf{das Bargeld} — \textquotedblleft{}the cash\textquotedblright{}.

\begin{quote}
Auf dem Markt zahlt man oft mit Bargeld.
\end{quote}

\begin{quote}
\emph{At the market people often pay in cash.}
\end{quote}

No plural. \emph{Bar zahlen} — to pay cash.

\begin{cousinweb}
\emph{Bargeld} is built from \emph{bar} 'bare, ready', an inherited Germanic word that is a cousin of English \emph{bare} (EtymDB), and \emph{Geld} 'money'.
\end{cousinweb}

\begin{grammarlens}[title={What you've built}]
One more word for telling, explaining and describing at B1.
\end{grammarlens}

\subsection*{Your turn}
\begin{itemize}
\raggedright
  \item \emph{Say it:} \emph{das Bargeld}
  \item \emph{Say it:} \emph{das Bargeld}, once more
  \item \emph{Say it:} a sentence of your own with \emph{das Bargeld}
  \item \emph{Recall:} say the German for the cashier, then the German for to take the money, then say \emph{das Bargeld} again
\end{itemize}

\begin{tcolorbox}[breakable,colback=teal!4,colframe=teal!35!black,title={Before you move on}]
What does das Bargeld mean? (\textquotedblleft{}The cash\textquotedblright{}.) Say it once more.
\end{tcolorbox}
\section[das Kartenlesegerät]{das Kartenlesegerät — the card reader, the card machine}

\label{lesson:GE-C449-kartenlesegeraet}



\begin{quote}
Before the new one: say the German for to take the money, then the German for the cash.
\end{quote}

\subsection*{You'll want to know: das Kartenlesegerät}
\textbf{das Kartenlesegerät} — \textquotedblleft{}the card reader, the card machine\textquotedblright{}.

\begin{quote}
Das Kartenlesegerät funktioniert heute nicht.
\end{quote}

\begin{quote}
\emph{The card machine is not working today.}
\end{quote}

Plural: \emph{die Kartenlesegeräte}.

\begin{cousinweb}
\emph{Kartenlesegerät} is built from \emph{Karte} 'card', from Latin \emph{charta} (Lewis \& Short), \emph{lesen} 'to read', inherited from Proto-Germanic \emph{lesaną}, and \emph{Gerät} 'device'.
\end{cousinweb}

\begin{grammarlens}[title={What you've built}]
One more word for telling, explaining and describing at B1.
\end{grammarlens}

\subsection*{Your turn}
\begin{itemize}
\raggedright
  \item \emph{Say it:} \emph{das Kartenlesegerät}
  \item \emph{Say it:} \emph{das Kartenlesegerät}, once more
  \item \emph{Say it:} a sentence of your own with \emph{das Kartenlesegerät}
  \item \emph{Recall:} say the German for to take the money, then the German for the cash, then say \emph{das Kartenlesegerät} again
\end{itemize}

\begin{tcolorbox}[breakable,colback=teal!4,colframe=teal!35!black,title={Before you move on}]
What does das Kartenlesegerät mean? (\textquotedblleft{}The card reader, the card machine\textquotedblright{}.) Say it once more.
\end{tcolorbox}
\section[das Wechselgeld]{das Wechselgeld — the change (money given back)}

\label{lesson:GE-C449-wechselgeld}



\begin{quote}
Before the new one: say the German for the cash, then the German for the card reader.
\end{quote}

\subsection*{You'll want to know: das Wechselgeld}
\textbf{das Wechselgeld} — \textquotedblleft{}the change (money given back)\textquotedblright{}.

\begin{quote}
Sie zählt das Wechselgeld und lächelt.
\end{quote}

\begin{quote}
\emph{She counts her change and smiles.}
\end{quote}

No plural. Small coins in a pocket are \emph{das Kleingeld}.

\begin{cousinweb}
\emph{Wechselgeld} is built from \emph{Wechsel} 'change, exchange', inherited from Old High German \emph{wehsal} (EtymDB), and \emph{Geld} 'money'.
\end{cousinweb}

\begin{grammarlens}[title={What you've built}]
One more word for telling, explaining and describing at B1.
\end{grammarlens}

\subsection*{Your turn}
\begin{itemize}
\raggedright
  \item \emph{Say it:} \emph{das Wechselgeld}
  \item \emph{Say it:} \emph{das Wechselgeld}, once more
  \item \emph{Say it:} a sentence of your own with \emph{das Wechselgeld}
  \item \emph{Recall:} say the German for the cash, then the German for the card reader, then say \emph{das Wechselgeld} again
\end{itemize}

\begin{tcolorbox}[breakable,colback=teal!4,colframe=teal!35!black,title={Before you move on}]
What does das Wechselgeld mean? (\textquotedblleft{}The change (money given back)\textquotedblright{}.) Say it once more.
\end{tcolorbox}
